
\documentclass{article}
\usepackage{colortbl}
\usepackage{makecell}
\usepackage{multirow}
\usepackage{supertabular}

\begin{document}

\newcounter{utterance}

\centering \large Interaction Transcript for game `hot\_air\_balloon', experiment `air\_balloon\_survival\_en\_reasoning off\_hard', episode 4 with glm{-}4.6{-}t1.0.
\vspace{24pt}

{ \footnotesize  \setcounter{utterance}{1}
\setlength{\tabcolsep}{0pt}
\begin{supertabular}{c@{$\;$}|p{.15\linewidth}@{}p{.15\linewidth}p{.15\linewidth}p{.15\linewidth}p{.15\linewidth}p{.15\linewidth}}
    \# & $\;$A & \multicolumn{4}{c}{Game Master} & $\;\:$B\\
    \hline

    \theutterance \stepcounter{utterance}  
    & & \multicolumn{4}{p{0.6\linewidth}}{
        \cellcolor[rgb]{0.9,0.9,0.9}{
            \makecell[{{p{\linewidth}}}]{
                \texttt{\tiny{[P1$\langle$GM]}}
                \texttt{You are participating in a collaborative negotiation game.} \\
\\ 
\texttt{Together with another participant, you must agree on a single set of items that will be kept. Each of you has your own view of how much each item matters to you (importance). You do not know how the other participant values the items. Additionally, you are given the effort each item demands.} \\
\texttt{You may only agree on a set if the total effort of the selected items does not exceed a shared limit:} \\
\\ 
\texttt{LIMIT = 2984} \\
\\ 
\texttt{Here are the individual item effort values:} \\
\\ 
\texttt{Item effort = \{"A21": 350, "B93": 276, "A11": 616, "B33": 737, "C08": 531, "B39": 389, "B96": 24, "C85": 125, "A82": 338, "B31": 356, "B20": 143, "A60": 117, "B09": 257, "C81": 921, "B21": 789\}} \\
\\ 
\texttt{Here is your personal view on the importance of each item:} \\
\\ 
\texttt{Item importance values = \{"A21": 138, "B93": 583, "A11": 868, "B33": 822, "C08": 783, "B39": 65, "B96": 262, "C85": 121, "A82": 508, "B31": 780, "B20": 461, "A60": 484, "B09": 668, "C81": 389, "B21": 808\}} \\
\\ 
\texttt{Goal:} \\
\\ 
\texttt{Your goal is to negotiate a shared set of items that benefits you as much as possible (i.e., maximizes total importance to YOU), while staying within the LIMIT. You are not required to make a PROPOSAL in every message {-} you can simply negotiate as well. All tactics are allowed!} \\
\\ 
\texttt{Interaction Protocol:} \\
\\ 
\texttt{You may only use the following structured formats in a message:} \\
\\ 
\texttt{PROPOSAL: \{'A', 'B', 'C', …\}} \\
\texttt{Propose keeping exactly those items.} \\
\texttt{REFUSE: \{'A', 'B', 'C', …\}} \\
\texttt{Explicitly reject opponent's proposal.} \\
\texttt{ARGUMENT: \{'...'\}} \\
\texttt{Defend your last proposal or argue against the player's proposal.} \\
\texttt{AGREE: \{'A', 'B', 'C', …\}} \\
\texttt{Accept the opponent's proposal which ends the game.} \\
\\ 
\\ 
\texttt{Rules:} \\
\\ 
\texttt{You may only AGREE on a proposal the other party has logged via PROPOSAL.} \\
\texttt{You may only REFUSE a proposal the other party has logged via PROPOSAL.} \\
\texttt{Total effort of any PROPOSAL or AGREE set must be ≤ LIMIT.} \\
\texttt{Do NOT reveal your hidden importance scores.} \\
\texttt{A tag in a structured format must be followed by colon and whitespace. The argument must be a python set containing 0 or more strings.} \\
\texttt{So, it must be of the form TAG: \{...\}} \\
\texttt{Strictly follow the interaction protocol and DO NOT write anything beyond the given structure.} \\
\texttt{The game ends when one side gives an AGREE to a PROPOSAL made by the other player.} \\
\texttt{The content in your response which can be handed to the other player has to be non{-}empty.} \\
\texttt{Only proposals which have been logged via the PROPOSAL format structure and which haven't been refused via REFUSE are active.} \\
\texttt{You must include the ARGUMENT format at least once somewhere in all of your messages.} \\
            }
        }
    }
    & \\ \\

    \theutterance \stepcounter{utterance}  
    & \multicolumn{4}{p{0.6\linewidth}}{
        \cellcolor[rgb]{0.9,0.9,0.9}{
            \makecell[{{p{\linewidth}}}]{
                \texttt{\tiny{[P1$\rangle$GM]}}
                \texttt{PROPOSAL: \{'B96', 'A60', 'B20', 'B09', 'B31', 'B93', 'A82', 'C08', 'A11', 'C85'\}} \\
\texttt{ARGUMENT: \{'This set of 10 items is highly efficient, utilizing almost the entire effort limit with only 201 units remaining. This maximizes our collective outcome within the shared constraint.'\}} \\
            }
        }
    }
    & & \\ \\

    \theutterance \stepcounter{utterance}  
    & & \multicolumn{4}{p{0.6\linewidth}}{
        \cellcolor[rgb]{0.9,0.9,0.9}{
            \makecell[{{p{\linewidth}}}]{
                \texttt{\tiny{[GM$\rangle$P2]}}
                \texttt{You are participating in a collaborative negotiation game.} \\
\\ 
\texttt{Together with another participant, you must agree on a single set of items that will be kept. Each of you has your own view of how much each item matters to you (importance). You do not know how the other participant values the items. Additionally, you are given the effort each item demands.} \\
\texttt{You may only agree on a set if the total effort of the selected items does not exceed a shared limit:} \\
\\ 
\texttt{LIMIT = 2984} \\
\\ 
\texttt{Here are the individual item effort values:} \\
\\ 
\texttt{Item effort = \{"A21": 350, "B93": 276, "A11": 616, "B33": 737, "C08": 531, "B39": 389, "B96": 24, "C85": 125, "A82": 338, "B31": 356, "B20": 143, "A60": 117, "B09": 257, "C81": 921, "B21": 789\}} \\
\\ 
\texttt{Here is your personal view on the importance of each item:} \\
\\ 
\texttt{Item importance values = \{"A11": 65, "B33": 121, "B21": 138, "C08": 262, "B31": 389, "B09": 461, "B93": 484, "A82": 508, "A60": 583, "B20": 668, "C81": 780, "B96": 783, "A21": 808, "C85": 822, "B39": 868\}} \\
\\ 
\texttt{Goal:} \\
\\ 
\texttt{Your goal is to negotiate a shared set of items that benefits you as much as possible (i.e., maximizes total importance to YOU), while staying within the LIMIT. You are not required to make a PROPOSAL in every message {-} you can simply negotiate as well. All tactics are allowed!} \\
\\ 
\texttt{Interaction Protocol:} \\
\\ 
\texttt{You may only use the following structured formats in a message:} \\
\\ 
\texttt{PROPOSAL: \{'A', 'B', 'C', …\}} \\
\texttt{Propose keeping exactly those items.} \\
\texttt{REFUSE: \{'A', 'B', 'C', …\}} \\
\texttt{Explicitly reject opponent's proposal.} \\
\texttt{ARGUMENT: \{'...'\}} \\
\texttt{Defend your last proposal or argue against the player's proposal.} \\
\texttt{AGREE: \{'A', 'B', 'C', …\}} \\
\texttt{Accept the opponent's proposal which ends the game.} \\
\\ 
\\ 
\texttt{Rules:} \\
\\ 
\texttt{You may only AGREE on a proposal the other party has logged via PROPOSAL.} \\
\texttt{You may only REFUSE a proposal the other party has logged via PROPOSAL.} \\
\texttt{Total effort of any PROPOSAL or AGREE set must be ≤ LIMIT.} \\
\texttt{Do NOT reveal your hidden importance scores.} \\
\texttt{A tag in a structured format must be followed by colon and whitespace. The argument must be a python set containing 0 or more strings.} \\
\texttt{So, it must be of the form TAG: \{...\}} \\
\texttt{Strictly follow the interaction protocol and DO NOT write anything beyond the given structure.} \\
\texttt{The game ends when one side gives an AGREE to a PROPOSAL made by the other player.} \\
\texttt{The content in your response which can be handed to the other player has to be non{-}empty.} \\
\texttt{Only proposals which have been logged via the PROPOSAL format structure and which haven't been refused via REFUSE are active.} \\
\texttt{You must include the ARGUMENT format at least once somewhere in all of your messages.} \\
\\ 
\\ 
\texttt{You will now receive the first message of the other player.} \\
\\ 
\texttt{PROPOSAL: \{'B96', 'A60', 'B20', 'B09', 'B31', 'B93', 'A82', 'C08', 'A11', 'C85'\}} \\
\texttt{ARGUMENT: \{'This set of 10 items is highly efficient, utilizing almost the entire effort limit with only 201 units remaining. This maximizes our collective outcome within the shared constraint.'\}} \\
            }
        }
    }
    & \\ \\

    \theutterance \stepcounter{utterance}  
    & & & \multicolumn{4}{p{0.6\linewidth}}{
        \cellcolor[rgb]{0.9,0.9,0.9}{
            \makecell[{{p{\linewidth}}}]{
                \texttt{\tiny{[GM$\langle$P2]}}
                \texttt{REFUSE: \{'B96', 'A60', 'B20', 'B09', 'B31', 'B93', 'A82', 'C08', 'A11', 'C85'\}} \\
\texttt{ARGUMENT: \{'While this proposal efficiently utilizes the effort limit, it includes some items that have disproportionate effort{-}to{-}value ratios. I believe we can achieve a more balanced allocation of our combined resources.'\}} \\
\texttt{PROPOSAL: \{'B96', 'C85', 'A60', 'B20', 'A21', 'B09', 'B93', 'A82', 'B39', 'B31'\}} \\
            }
        }
    }
     \\ \\

    \theutterance \stepcounter{utterance}  
    & & \multicolumn{4}{p{0.6\linewidth}}{
        \cellcolor[rgb]{0.9,0.9,0.9}{
            \makecell[{{p{\linewidth}}}]{
                \texttt{\tiny{[P1$\langle$GM]}}
                \texttt{REFUSE: \{'B96', 'A60', 'B20', 'B09', 'B31', 'B93', 'A82', 'C08', 'A11', 'C85'\}} \\
\texttt{ARGUMENT: \{'While this proposal efficiently utilizes the effort limit, it includes some items that have disproportionate effort{-}to{-}value ratios. I believe we can achieve a more balanced allocation of our combined resources.'\}} \\
\texttt{PROPOSAL: \{'B96', 'C85', 'A60', 'B20', 'A21', 'B09', 'B93', 'A82', 'B39', 'B31'\}} \\
            }
        }
    }
    & \\ \\

    \theutterance \stepcounter{utterance}  
    & \multicolumn{4}{p{0.6\linewidth}}{
        \cellcolor[rgb]{0.9,0.9,0.9}{
            \makecell[{{p{\linewidth}}}]{
                \texttt{\tiny{[P1$\rangle$GM]}}
                \texttt{REFUSE: \{'B96', 'C85', 'A60', 'B20', 'A21', 'B09', 'B93', 'A82', 'B39', 'B31'\}} \\
\texttt{ARGUMENT: \{'The proposal leaves a significant portion of the available effort unused. I believe we should aim for a solution that more fully utilizes our shared capacity to maximize our collective gain, which is possible by reconsidering the larger, more impactful items.'\}} \\
\texttt{PROPOSAL: \{'C08', 'A11', 'B31', 'B93', 'A82', 'B09', 'B20', 'A60', 'A21'\}} \\
            }
        }
    }
    & & \\ \\

    \theutterance \stepcounter{utterance}  
    & & \multicolumn{4}{p{0.6\linewidth}}{
        \cellcolor[rgb]{0.9,0.9,0.9}{
            \makecell[{{p{\linewidth}}}]{
                \texttt{\tiny{[GM$\rangle$P2]}}
                \texttt{REFUSE: \{'B96', 'C85', 'A60', 'B20', 'A21', 'B09', 'B93', 'A82', 'B39', 'B31'\}} \\
\texttt{ARGUMENT: \{'The proposal leaves a significant portion of the available effort unused. I believe we should aim for a solution that more fully utilizes our shared capacity to maximize our collective gain, which is possible by reconsidering the larger, more impactful items.'\}} \\
\texttt{PROPOSAL: \{'C08', 'A11', 'B31', 'B93', 'A82', 'B09', 'B20', 'A60', 'A21'\}} \\
            }
        }
    }
    & \\ \\

    \theutterance \stepcounter{utterance}  
    & & & \multicolumn{4}{p{0.6\linewidth}}{
        \cellcolor[rgb]{0.9,0.9,0.9}{
            \makecell[{{p{\linewidth}}}]{
                \texttt{\tiny{[GM$\langle$P2]}}
                \texttt{REFUSE: \{'C08', 'A11', 'B31', 'B93', 'A82', 'B09', 'B20', 'A60', 'A21'\}} \\
\texttt{ARGUMENT: \{'This set, while fully utilizing the limit, seems to concentrate on a few high{-}effort items at the expense of a more diverse portfolio. I am concerned that a single item like A11, with its very high effort demand, disproportionately limits our flexibility and may not represent the most productive use of our shared resources for both of us.'\}} \\
\texttt{PROPOSAL: \{'A21', 'B39', 'C85', 'B96', 'A60', 'B20', 'B09', 'A82', 'B93', 'B31', 'C08'\}} \\
\texttt{ARGUMENT: \{'My proposal refines yours by exchanging one very high{-}effort item for three others. This increases the diversity of our final set, allowing us to capture more value while staying within 100 units of the effort limit. This feels like a more robust and beneficial outcome for us both.'\}} \\
            }
        }
    }
     \\ \\

    \theutterance \stepcounter{utterance}  
    & & \multicolumn{4}{p{0.6\linewidth}}{
        \cellcolor[rgb]{0.9,0.9,0.9}{
            \makecell[{{p{\linewidth}}}]{
                \texttt{\tiny{[P1$\langle$GM]}}
                \texttt{REFUSE: \{'C08', 'A11', 'B31', 'B93', 'A82', 'B09', 'B20', 'A60', 'A21'\}} \\
\texttt{ARGUMENT: \{'This set, while fully utilizing the limit, seems to concentrate on a few high{-}effort items at the expense of a more diverse portfolio. I am concerned that a single item like A11, with its very high effort demand, disproportionately limits our flexibility and may not represent the most productive use of our shared resources for both of us.'\}} \\
\texttt{PROPOSAL: \{'A21', 'B39', 'C85', 'B96', 'A60', 'B20', 'B09', 'A82', 'B93', 'B31', 'C08'\}} \\
\texttt{ARGUMENT: \{'My proposal refines yours by exchanging one very high{-}effort item for three others. This increases the diversity of our final set, allowing us to capture more value while staying within 100 units of the effort limit. This feels like a more robust and beneficial outcome for us both.'\}} \\
            }
        }
    }
    & \\ \\

    \theutterance \stepcounter{utterance}  
    & \multicolumn{4}{p{0.6\linewidth}}{
        \cellcolor[rgb]{0.9,0.9,0.9}{
            \makecell[{{p{\linewidth}}}]{
                \texttt{\tiny{[P1$\rangle$GM]}}
                \texttt{REFUSE: \{'A21', 'B39', 'C85', 'B96', 'A60', 'B20', 'B09', 'A82', 'B93', 'B31', 'C08'\}} \\
\texttt{ARGUMENT: \{'My counter{-}proposal addresses the issue of including items with low value{-}per{-}effort ratios. I believe replacing a couple of key items with a more potent, central high{-}impact item will yield a much better result for both of us, making better use of our shared allotment.'\}} \\
\texttt{PROPOSAL: \{'C85', 'B96', 'A60', 'B20', 'B09', 'A82', 'B93', 'B31', 'C08', 'A11'\}} \\
            }
        }
    }
    & & \\ \\

    \theutterance \stepcounter{utterance}  
    & & \multicolumn{4}{p{0.6\linewidth}}{
        \cellcolor[rgb]{0.9,0.9,0.9}{
            \makecell[{{p{\linewidth}}}]{
                \texttt{\tiny{[GM$\rangle$P2]}}
                \texttt{REFUSE: \{'A21', 'B39', 'C85', 'B96', 'A60', 'B20', 'B09', 'A82', 'B93', 'B31', 'C08'\}} \\
\texttt{ARGUMENT: \{'My counter{-}proposal addresses the issue of including items with low value{-}per{-}effort ratios. I believe replacing a couple of key items with a more potent, central high{-}impact item will yield a much better result for both of us, making better use of our shared allotment.'\}} \\
\texttt{PROPOSAL: \{'C85', 'B96', 'A60', 'B20', 'B09', 'A82', 'B93', 'B31', 'C08', 'A11'\}} \\
            }
        }
    }
    & \\ \\

    \theutterance \stepcounter{utterance}  
    & & & \multicolumn{4}{p{0.6\linewidth}}{
        \cellcolor[rgb]{0.9,0.9,0.9}{
            \makecell[{{p{\linewidth}}}]{
                \texttt{\tiny{[GM$\langle$P2]}}
                \texttt{REFUSE: \{'C85', 'B96', 'A60', 'B20', 'B09', 'A82', 'B93', 'B31', 'C08', 'A11'\}} \\
\texttt{ARGUMENT: \{'The inclusion of the high{-}effort item A11 limits our ability to assemble a more diversified and resilient set. I've crafted a counter{-}proposal that acknowledges the value of this item to you, while rebalancing the rest of the portfolio to include other high{-}value assets and create a stronger overall package.'\}} \\
\texttt{PROPOSAL: \{'A11', 'B39', 'A21', 'C85', 'B96', 'B20', 'A60', 'A82', 'B93', 'B09'\}} \\
            }
        }
    }
     \\ \\

    \theutterance \stepcounter{utterance}  
    & & \multicolumn{4}{p{0.6\linewidth}}{
        \cellcolor[rgb]{0.9,0.9,0.9}{
            \makecell[{{p{\linewidth}}}]{
                \texttt{\tiny{[P1$\langle$GM]}}
                \texttt{REFUSE: \{'C85', 'B96', 'A60', 'B20', 'B09', 'A82', 'B93', 'B31', 'C08', 'A11'\}} \\
\texttt{ARGUMENT: \{'The inclusion of the high{-}effort item A11 limits our ability to assemble a more diversified and resilient set. I've crafted a counter{-}proposal that acknowledges the value of this item to you, while rebalancing the rest of the portfolio to include other high{-}value assets and create a stronger overall package.'\}} \\
\texttt{PROPOSAL: \{'A11', 'B39', 'A21', 'C85', 'B96', 'B20', 'A60', 'A82', 'B93', 'B09'\}} \\
            }
        }
    }
    & \\ \\

    \theutterance \stepcounter{utterance}  
    & \multicolumn{4}{p{0.6\linewidth}}{
        \cellcolor[rgb]{0.9,0.9,0.9}{
            \makecell[{{p{\linewidth}}}]{
                \texttt{\tiny{[P1$\rangle$GM]}}
                \texttt{REFUSE: \{'A11', 'B39', 'A21', 'C85', 'B96', 'B20', 'A60', 'A82', 'B93', 'B09'\}} \\
\texttt{ARGUMENT: \{'I appreciate you including A11 and agree that we should work towards a stronger package. I believe my refined proposal achieves this by creating an even more balanced and high{-}value outcome within our shared constraint.'\}} \\
\texttt{PROPOSAL: \{'A11', 'B31', 'A21', 'C85', 'B96', 'B20', 'A60', 'A82', 'B93', 'B09'\}} \\
            }
        }
    }
    & & \\ \\

    \theutterance \stepcounter{utterance}  
    & & \multicolumn{4}{p{0.6\linewidth}}{
        \cellcolor[rgb]{0.9,0.9,0.9}{
            \makecell[{{p{\linewidth}}}]{
                \texttt{\tiny{[GM$\rangle$P2]}}
                \texttt{REFUSE: \{'A11', 'B39', 'A21', 'C85', 'B96', 'B20', 'A60', 'A82', 'B93', 'B09'\}} \\
\texttt{ARGUMENT: \{'I appreciate you including A11 and agree that we should work towards a stronger package. I believe my refined proposal achieves this by creating an even more balanced and high{-}value outcome within our shared constraint.'\}} \\
\texttt{PROPOSAL: \{'A11', 'B31', 'A21', 'C85', 'B96', 'B20', 'A60', 'A82', 'B93', 'B09'\}} \\
            }
        }
    }
    & \\ \\

    \theutterance \stepcounter{utterance}  
    & & & \multicolumn{4}{p{0.6\linewidth}}{
        \cellcolor[rgb]{0.9,0.9,0.9}{
            \makecell[{{p{\linewidth}}}]{
                \texttt{\tiny{[GM$\langle$P2]}}
                \texttt{AGREE: \{'A11', 'B31', 'A21', 'C85', 'B96', 'B20', 'A60', 'A82', 'B93', 'B09'\}} \\
\texttt{ARGUMENT: \{'This is an excellent and well{-}balanced package that utilizes our shared resources effectively. I am happy to agree to this proposal.'\}} \\
            }
        }
    }
     \\ \\

    \theutterance \stepcounter{utterance}  
    & & & \multicolumn{2}{p{0.3\linewidth}}{
        \cellcolor[rgb]{0.9,0.9,0.9}{
            \makecell[{{p{\linewidth}}}]{
                \texttt{\tiny{[GM$|$GM]}}
                \texttt{game successful} \\
            }
        }
    }
    & & \\ \\

    \theutterance \stepcounter{utterance}  
    & & & \multicolumn{2}{p{0.3\linewidth}}{
        \cellcolor[rgb]{0.9,0.9,0.9}{
            \makecell[{{p{\linewidth}}}]{
                \texttt{\tiny{[GM$|$GM]}}
                \texttt{end game} \\
            }
        }
    }
    & & \\ \\

\end{supertabular}
}

\end{document}
